% =====================================================================
%  DLoc Compression — Defense Q&A (Results & Discussion -> Conclusion)
%  Self-contained: compile with pdfLaTeX (Overleaf-ready).
% =====================================================================
\documentclass[11pt]{article}

\usepackage[a4paper,margin=2.2cm]{geometry}
\usepackage{amsmath,amssymb}
\usepackage{booktabs}
\usepackage{xcolor}
\usepackage{enumitem}
\usepackage{tcolorbox}
\tcbuselibrary{skins,breakable}
\usepackage[hidelinks]{hyperref}

\newtcolorbox{trap}[1][TRAP --- be careful]{%
  colback=red!4, colframe=red!55!black, fonttitle=\bfseries, title=#1,
  breakable, left=2mm, right=2mm, top=1mm, bottom=1mm}
\newtcolorbox{safe}[1][Safe phrasing]{%
  colback=green!5, colframe=green!45!black, fonttitle=\bfseries, title=#1,
  breakable, left=2mm, right=2mm, top=1mm, bottom=1mm}

% Q/A formatting
\newcommand{\Q}[1]{\medskip\noindent\textbf{\textcolor{blue!50!black}{Q.}}\ \textbf{#1}\par\nopagebreak}
\newcommand{\A}{\noindent\textbf{A.}\ }

\title{\textbf{DLoc Compression --- Defense Q\&A}\\[2mm]
\large Results \& Discussion, Findings, and Conclusion}
\date{}
\begin{document}
\maketitle
\tableofcontents
\vspace{3mm}\hrule\vspace{3mm}

\noindent This is an anticipated-questions bank for the defense. Numbers are exact. Red boxes mark
\textbf{traps} --- questions where a loose answer can get you in trouble --- with the safe phrasing.

% =====================================================================
\section{In-Distribution Accuracy (Experiment 1)}
% =====================================================================

\Q{What does ``in-distribution'' mean, and what is the test set?}
\A Train and test in the \emph{same} environment (July). We train on 17{,}574 July samples and test on a
\textbf{held-out} 4{,}260-sample split of Session~2 (collected $\sim$1\,h later). It measures best-case
accuracy --- ``can the model learn the mapping when nothing changes?''

\Q{Is your in-distribution test set just the training data?}
\A No. It is a \textbf{held-out} split never seen during training. Same \emph{environment} (so it is
``in-distribution''), but genuinely held-out data, not memorized samples.

\Q{Summarize the in-distribution results.}
\A MobileNetV2 \textbf{matches} DLoc at 0.707\,m (7.2$\times$ smaller); MobileNetV2-UNet is best at
\textbf{0.640\,m} (6.8$\times$ smaller); TinyCNN fails at 3.92\,m; Mamba fails at 10.56\,m.

\Q{Why does MobileNetV2 match DLoc with 7$\times$ fewer parameters?}
\A Because DLoc is \textbf{over-parameterized} for in-distribution modelling: 2.3M parameters are enough
to learn the CSI$\to$location mapping, and depthwise-separable convolutions learn it efficiently.

\Q{Why does the U-Net beat the plain MobileNetV2?}
\A The \textbf{skip connections} carry fine, high-resolution spatial detail past the bottleneck, so the
predicted bump is sharper --- 0.640\,m vs 0.707\,m.

\Q{Why does TinyCNN fail?}
\A It has only 47K parameters --- below the $\sim$2M \textbf{capacity floor}. With three conv layers it
cannot represent both the per-AP features and the spatial mapping. No amount of training fixes too little
capacity.

\Q{Why does Mamba fail so badly (10.56\,m)?}
\A \textbf{Architecture mismatch}, not lack of capacity. Mamba treats the input as a 1D sequence, but the
target is a 2D spatial heatmap; flattening destroys spatial locality. It reaches its best at epoch~1 and
degrades monotonically (train and test error rise together) --- which \emph{rules out overfitting}.

\Q{Why didn't distillation improve in-distribution accuracy?}
\A In-distribution, the teacher's output is nearly identical to the ground truth, so its supervision is
\textbf{redundant}. Distillation helps only \emph{outside} the training distribution (noise,
cross-session).

% =====================================================================
\section{U-Net Architecture and Scheduler (Experiment 1b)}
% =====================================================================

\Q{What is the difference between cosine and step learning-rate schedules?}
\A Cosine \textbf{smoothly} anneals the LR from $10^{-4}$ to $\sim$0 over 120 epochs; step \textbf{halves}
the LR every 20 epochs (a staircase). Both reach median 0.640\,m.

\Q{Which scheduler is better, and by how much?}
\A Cosine, marginally: same median (0.640), slightly tighter tails (P90 1.523 / P99 2.750 vs step's
1.552 / 2.918). Its smooth decay toward zero in the last epochs fine-tunes the worst cases better. The
\textbf{architecture dominates}; the scheduler has a small effect.

\Q{Which exact model is the green ``U-Net'' curve in the CDF figure?}
\A The \textbf{cosine-annealing, standalone} MobileNetV2-UNet (\texttt{unet\_cosine\_standalone\_best.pth}
in \texttt{dump\_results.py}).

% =====================================================================
\section{Robustness (Experiment 2)}
% =====================================================================

\Q{How is robustness evaluated --- did you retrain?}
\A No retraining. We take the in-distribution weights and apply corruptions to the test inputs: AP
dropout, Gaussian noise ($\sigma{=}0.1,0.2$), attenuation (0.5$\times$ global, 0.3$\times$ one AP), and
5$\times$5 Gaussian blur.

\Q{What are the two headline robustness findings?}
\A \textbf{Distillation $\to$ 3.4$\times$} noise robustness (MobileNetV2 at $\sigma{=}0.1$: 16.4 $\to$
4.8\,m); \textbf{skip connections $\to$ 5.5$\times$} (standalone U-Net 3.0\,m vs MobileNetV2 16.4\,m).

\Q{Why does distillation make a model more noise-robust?}
\A The teacher's softer heatmaps act as \textbf{spatial label smoothing}; the student learns a smoother
representation that noise perturbs less --- the spatial-regression analogue of soft targets in
classification distillation.

\Q{Why do skip connections give robustness?}
\A When noise corrupts the deep bottleneck, the \textbf{clean shallow features bypass it} through the
skips --- a second, uncorrupted information path the model can fall back on.

\Q{Why does distillation \emph{harm} TinyCNN?}
\A Too little capacity. TinyCNN spends its 47K parameters learning the teacher's model of AP
interactions; when an AP drops at test time those assumptions become invalid and it has \textbf{no spare
capacity to fall back on} (1-AP dropout: 7.9 $\to$ 13.6\,m). Hence the $\sim$2M capacity threshold.

\Q{How did you discover the $\sim$2M capacity threshold?}
\A By toggling distillation on/off across model sizes and watching its effect \textbf{flip sign}: it
helped the $\geq$2M students (MobileNetV2 2.3M, U-Net 2.41M) but harmed 47K TinyCNN. The boundary lies
between them, so we state ``roughly 2 million parameters.''

\Q{Why ``complementary but not additive''?}
\A Skips and distillation fix robustness by \emph{different} mechanisms (bypassing the bottleneck vs
smoothing the representation), so they don't simply add. On the U-Net under noise, the standalone (skips
only) is 3.0\,m and adding distillation gives 3.47\,m --- the skip robustness already dominates. Skips
protect against high noise; distillation protects against multi-AP dropout.

\Q{Which corruption is worst, and why?}
\A \textbf{Gaussian noise} --- it \emph{adds fake structure} (false peaks the model can be fooled by),
sending fragile models to 16\,m. Attenuation and blur only \emph{weaken or soften} the real signal (the
true peak stays put), so they cause far smaller errors (blur barely moves the U-Net: 0.67\,m).

% =====================================================================
\section{Cross-Session Generalization (Experiment 3)}
% =====================================================================

\Q{What is leave-one-session-out (LOSO), and why use it?}
\A Three folds; each trains on July + two August sessions and tests on the \textbf{held-out} August
session. We use it because it is the \emph{exact protocol the DLoc paper used}, so the comparison is fair.

\Q{Why train the U-Net standalone (not distilled) for cross-session?}
\A Distillation needs a teacher trained on the same data; to distill per fold we would have to retrain a
16.5M DLoc teacher each time --- expensive and it would make the DLoc comparison unfair. Standalone is
what is directly comparable to DLoc's published numbers.

\Q{Give the cross-session results.}
\A Median 0.852 / 0.822 / 1.116\,m on Aug16\_1 / Aug16\_3 / Aug16\_4\_ref (mean \textbf{0.930\,m}) vs
DLoc's 0.71 / 0.82 / 1.05 (mean 0.86), at 6.8$\times$ fewer parameters. We \textbf{match} DLoc on the
different-furniture fold and have a \emph{lower} P90 on the reflector fold.

\Q{Why is the reflector fold the hardest?}
\A An \textbf{aluminum reflector} is an almost ideal WiFi mirror; it injects strong specular ``phantom''
reflections (extra fake blobs), the most severe input distortion.

\Q{Why is DLoc slightly better on the reflector fold ($\sim$0.07\,m)?}
\A DLoc's \textbf{consistency decoder} enforces agreement across APs, which helps most under strong
specular multipath; our single-decoder U-Net has no such mechanism. The gap is small and expected.

\Q{How do you know the cross-session result isn't luck?}
\A Three random seeds (42/123/777) per fold; the median standard deviation is only 0.016--0.044\,m ---
stable.

\Q{Why use median error rather than mean?}
\A The median is robust to a few wild outliers (ambiguous locations) and is the metric the DLoc paper
reports, so it keeps the comparison like-for-like. We also report P90/P99 and CDFs for the tails.

% =====================================================================
\section{Zero-Shot Cross-Session (Experiment 4)}
% =====================================================================

\Q{What is the zero-shot evaluation, and why isn't it the DLoc comparison?}
\A Train on July \emph{only}, test on the August sessions (never seen). It is the harshest domain shift,
used as an internal \textbf{ablation} to compare architectures/distillation --- deliberately \emph{not} a
fair DLoc comparison.

\Q{What did zero-shot show?}
\A Error rises across the August sessions and is highest on the reflector for \emph{every} model --- a
difficulty hierarchy that is a property of the \textbf{environment}, not the model. Distillation helped
only on the hardest (reflector) fold (e.g.\ U-Net $1.334 \to 1.281$\,m).

% =====================================================================
\section{Multi-Seed Validation and Statistics (Experiment 5)}
% =====================================================================

\Q{Why train the baseline with three seeds?}
\A To show the student-vs-baseline gap is not due to random initialization --- to characterize the
baseline's seed-to-seed variance.

\Q{What is the multi-seed baseline result?}
\A $0.875 \pm 0.036$\,m (seeds 0.825 / 0.906 / 0.894), over 120 epochs.

\Q{Why is the multi-seed baseline (0.875) higher than the single run (0.707)?}
\A The multi-seed runs use the full 120-epoch schedule; because StepLR ($\gamma{=}0.9$) \textbf{never
decays the LR to zero}, longer training drifts slightly upward. The 0.707\,m was a favorable 50-epoch run.

\Q{Is the U-Net's improvement statistically significant?}
\A Yes. The U-Net's 0.640\,m is \textbf{6.5 standard deviations} below the multi-seed baseline mean
(0.875\,m, std 0.036) --- far larger than seed noise.

\Q{Why does P99 vary so much across seeds ($\pm6.8$\,m)?}
\A P99 depends on a \emph{few} inherently ambiguous locations whose accuracy varies between seeds; with
so few samples driving it, it is a noisy statistic (CV 52\%), unlike the stable median (CV 4.1\%).

% =====================================================================
\section{Computational Efficiency (Experiment 6)}
% =====================================================================

\Q{Quantify the compression.}
\A U-Net: 2.41M params (6.8$\times$), \textbf{1.45\,GFLOPs} (56$\times$ less than DLoc's 81.06\,G), and
\textbf{27.2\,ms} CPU (10$\times$ faster than 270.9\,ms); also 3.6$\times$ faster on GPU.

\Q{Why is the FLOP reduction (56$\times$) bigger than the parameter reduction (6.8$\times$)?}
\A \textbf{Depthwise-separable convolutions} replace DLoc's full convolutions; they cut the multiply-add
operations far more than they cut the raw parameter count.

\Q{Why report CPU latency specifically?}
\A CPU is the platform most representative of \textbf{edge deployment} (IoT gateways, APs, phones without
a GPU) --- the whole motivation for compressing DLoc.

\Q{Which model is the ``DLoc'' efficiency figure?}
\A Its test-time subnet (encoder + location decoder); the consistency/offset decoder is training-only and
excluded.

% =====================================================================
\section{Knowledge Distillation (concepts)}
% =====================================================================

\Q{Is the teacher the same as the baseline?}
\A Yes --- one trained DLoc model (16.5M), in two roles: \textbf{baseline} = the accuracy yardstick;
\textbf{teacher} = the output source for distillation.

\Q{Do you copy the teacher's weights into the student?}
\A No. We run the teacher \emph{once} to record its output heatmaps (cached, $\sim$4.09\,GB train /
0.99\,GB test via \texttt{precompute\_teacher.py}); the student learns its \textbf{own} weights by
matching both the ground truth and those cached outputs. The teacher's weights are never updated and are
not loaded during student training.

\Q{Write the distillation loss.}
\A $L = 0.7\,\mathrm{MSE}(\hat{y}, y_{\text{true}}) + 0.3\,\mathrm{MSE}(\hat{y}, y_{\text{teacher}})$,
i.e.\ $\alpha{=}0.7$, $\beta{=}0.3$. Standalone sets $\beta{=}0$.

\Q{Why is this distillation ``dual-natured''?}
\A The same soft-target signal \emph{helps} robustness/generalization for models with capacity but
\emph{harms} the smallest model (TinyCNN) --- two opposite effects depending on capacity.

\Q{How does the U-Net get a ``clean-signal'' benefit like DLoc's consistency decoder without having one?}
\A Both keep the clean signal alive past the noisy bottleneck. DLoc \emph{trains a second decoder} to
reconstruct clean heatmaps; the U-Net \emph{wires the clean shallow features straight through} with skip
connections --- same goal, achieved structurally for free (no extra decoder, fewer parameters).

% =====================================================================
\section{Findings and Conclusion}
% =====================================================================

\Q{State the thesis's main conclusion in one sentence.}
\A A model roughly an order of magnitude smaller than DLoc achieves the same accuracy --- and
architectural (skip connections) and training-based (distillation) choices each buy measurable
robustness, with distillation only helping above a $\sim$2M capacity floor.

\Q{What is your practical recommendation?}
\A The distilled MobileNetV2-UNet is the best all-rounder (0.640\,m, competitive cross-session,
6.8$\times$ smaller, best under multi-AP dropout, 56$\times$ cheaper). If pure noise robustness matters,
the \emph{standalone} U-Net; if maximum compression matters, the distilled plain MobileNetV2 (0.707\,m,
2.3M).

\Q{What are the limitations?}
\A Single building (Jacobs Hall) --- no cross-building test; the median is pixel-quantized
($\sim0.1\sqrt{n}$), so results within $\sim$0.1\,m are effectively equal; and we compare to DLoc's
published numbers, not a re-run of DLoc under identical conditions.

\Q{What is the future work?}
\A Deploy on real edge hardware and collect a new multi-building dataset (e.g.\ at E-JUST) to test
cross-building generalization.

% =====================================================================
\section{Trap Questions (answer carefully)}
% =====================================================================

\begin{trap}[``You beat the published DLoc'' --- did you?]
The DLoc paper reports \textbf{64\,cm} (complex space, the setup you compare to) and 65\,cm headline.
Your U-Net's 0.640\,m \textbf{ties} the published 0.64\,m on median (and is slightly better on the tail:
P90 1.52 vs 1.6, P99 2.75 vs 3.2). Your ``9.5\% improvement'' is vs your \emph{own} re-implementation
(0.707\,m), \emph{not} the published number.
\end{trap}
\begin{safe}
``Our U-Net \textbf{matches DLoc's published 0.64\,m} median (with a slightly tighter tail) at
6.8$\times$ fewer parameters, and improves \textbf{9.5\% over our own faithful re-implementation}
(0.707\,m).''
\end{safe}

\begin{trap}[``Is 9.5\% the seed-vs-seed comparison?'']
No. 9.5\% is U-Net 0.640 vs the \emph{single} baseline run 0.707. There are \textbf{no in-distribution
U-Net seeds} (multi-seed was run for the baseline only; the U-Net's 3 seeds are for cross-session). The
seed comparison is the separate \textbf{6.5$\sigma$} statement against the multi-seed baseline (0.875).
\end{trap}

\begin{trap}[``Your medians are identical (0.640) across U-Net variants --- is that data leakage?'']
No --- it is \textbf{metric quantization}. The error lives on a 0.1\,m pixel grid, so the median snaps to
a $\sim0.1\sqrt{n}$ lattice; near-identical models land on the same lattice point. That is why we also
report mean, P90/P99, and CDFs, which \emph{do} separate them.
\end{trap}

\begin{trap}[``Mamba just needed more training / tuning'']
No. Its best epoch is the \emph{first} one and train+test error rise together afterwards --- that is the
signature of an \textbf{architectural mismatch} (1D sequence vs 2D heatmap), not under-training or
overfitting. More epochs make it worse, not better.
\end{trap}

\begin{trap}[``In-distribution and cross-session --- which 0.6/0.9 number is which?'']
\textbf{0.640\,m = in-distribution} (same July session, optimistic). \textbf{0.930\,m = cross-session
LOSO} (held-out August environment, honest). Never mix them.
\end{trap}

\begin{trap}[``Why is your multi-seed baseline (0.875) worse than your headline baseline (0.707)?'']
Different schedule length: the multi-seed runs are 120 epochs and StepLR ($\gamma{=}0.9$) never decays the
LR to zero, so longer training drifts up. Both are honestly reported; the U-Net (0.640) beats \emph{both}.
\end{trap}

\end{document}
